% Build in this directory: pdflatex -interaction=nonstopmode -halt-on-error resume2PT26.tex
% Style follows ../class1/class1PT26.tex; no guidance or syllabus page.
\documentclass[a4paper,10pt,pdflatex,ja=standard,jafont=ipaex-type1,
  magstyle=real,scale=1]{bxjsarticle}
\usepackage{lmodern}
\usepackage{amsmath,amssymb,mathtools}
\usepackage{geometry}
\geometry{reset,left=22mm,right=22mm,top=20mm,bottom=17mm,
  headheight=13pt,headsep=6mm,footskip=10mm}
\usepackage{xcolor}
\usepackage{enumitem}
\usepackage{array,booktabs,tabularx}
\usepackage{tikz}
\usepackage{fancyhdr}
\usepackage{hyperref}

\definecolor{ink}{HTML}{193A5C}
\definecolor{accent}{HTML}{237D82}
\definecolor{soft}{HTML}{EAF3F3}
\definecolor{muted}{HTML}{52616B}
\newcommand{\coursename}{確率論と統計学 I}
\hypersetup{unicode=true,colorlinks=true,urlcolor=accent,linkcolor=ink,
  pdfauthor={中島秀太},pdftitle={第2回：確率変数と確率分布}}
\pagestyle{fancy}
\fancyhf{}
\fancyhead[L]{\small\color{ink}\coursename{}・授業要約と演習}
\fancyhead[R]{\small\color{ink}第2回}
\fancyfoot[C]{\small\color{ink}\thepage/\pageref*{lastpage}}
\renewcommand{\headrulewidth}{0.3pt}
\setlength{\parindent}{0pt}
\setlength{\parskip}{2pt}
\setlength{\abovedisplayskip}{4pt}
\setlength{\belowdisplayskip}{4pt}
\setlength{\abovedisplayshortskip}{3pt}
\setlength{\belowdisplayshortskip}{4pt}
\setlist[itemize]{leftmargin=1.8em,itemsep=1pt,topsep=2pt,parsep=0pt}
\setlist[enumerate]{leftmargin=2em,itemsep=3pt,topsep=3pt,parsep=0pt}
\newcommand{\topic}[1]{\par\vspace{5pt}
  {\large\sffamily\mdseries\color{accent}#1}\par\vspace{2pt}}
\newcommand{\handouttitle}[1]{%
  \begin{center}
    {\Large\color{ink}#1}\par\vspace{3mm}
    {\color{accent}\rule{0.88\linewidth}{0.7pt}}
  \end{center}\vspace{-2mm}}
\newcommand{\keybox}[1]{\par\vspace{3pt}\begingroup
  \setlength{\fboxsep}{5pt}
  \colorbox{soft}{\parbox{\dimexpr\linewidth-2\fboxsep\relax}{#1}}%
  \endgroup\par\vspace{2pt}}
\newcommand{\examplelabel}[1]{{\sffamily\mdseries\color{ink}#1}}
\newcommand{\F}{\mathcal F}
\newcommand{\R}{\mathbb R}

\begin{document}
\fontsize{10pt}{13.5pt}\selectfont
\begingroup
\fancyhead[L]{\small\color{ink}\coursename{}・前回の演習解答}
\handouttitle{第1回\quad 演習問題の解答}

\topic{演習1：サイコロ2回の確率空間}
\keybox{公平なサイコロを独立に2回投げる試行の離散確率空間
$(\Omega,\F,P)$ を構成し、出た目の合計が3となる確率を求めよ。}

1回目の出目を $i$、2回目の出目を $j$ として、
$$
  \Omega=\{1,2,3,4,5,6\}^2
    =\{(i,j):1\le i,j\le6\},\qquad \F=2^\Omega
$$
とおく。公平性と独立性より、各結果の重みは
$$
  p_{(i,j)}=\frac16\cdot\frac16=\frac1{36},\qquad
  \sum_{(i,j)\in\Omega}p_{(i,j)}=36\cdot\frac1{36}=1.
$$
したがって、$|A|$ を集合 $A$ の要素数として、
$$
  P(A)=\sum_{(i,j)\in A}p_{(i,j)}=\frac{|A|}{36}
  \qquad(A\in\F)
$$
と定める。$P(\Omega)=1$ であり、互いに素な事象列 $(A_n)_{n\ge1}$ について
$$
  P\!\left(\bigcup_{n=1}^\infty A_n\right)
  =\frac1{36}\sum_{n=1}^\infty|A_n|
  =\sum_{n=1}^\infty P(A_n)
$$
が成り立つので、$(\Omega,\F,P)$ は求める離散確率空間である。
合計が3となる事象は $S=\{(i,j)\in\Omega:i+j=3\}=\{(1,2),(2,1)\}$ だから、
$$
  {P(S)=\frac2{36}=\frac1{18}}.
$$

\topic{演習2：共通部分の可測性と和事象の確率}
\keybox{$(\Omega,\F,P)$ を確率空間とする。任意の $A,B\in\F$ に対して、
$A\cap B\in\F$ および $P(A\cup B)=P(A)+P(B)-P(A\cap B)$ を示せよ。}

$A^c=\Omega\setminus A$ とする。$\sigma$-加法族は補集合と有限和について閉じているので、
ド・モルガンの法則より
$$
  A\cap B=(A^c\cup B^c)^c\in\F.
$$
また、$\varnothing=\Omega^c\in\F$ である。可算加法性に
$A_1=\Omega$, $A_n=\varnothing$ $(n\ge2)$ を代入すると、
$1=1+\sum_{n=2}^\infty P(\varnothing)$ より $P(\varnothing)=0$。
したがって、可算加法性で残りの事象を空集合とすれば有限加法性が得られる。

$$
  C=A\setminus B=A\cap B^c,\qquad
  D=B\setminus A=B\cap A^c,\qquad E=A\cap B
$$
とおく。いずれも $\F$ に属し、互いに素である。
$\sqcup$ は互いに素な集合の和を表すものとすると、
$$
  A=C\sqcup E,\qquad B=D\sqcup E,\qquad
  A\cup B=C\sqcup D\sqcup E.
$$
有限加法性から
$$
  \begin{aligned}
    P(A)+P(B)-P(A\cap B)
      &=\bigl(P(C)+P(E)\bigr)+\bigl(P(D)+P(E)\bigr)-P(E)\\
      &=P(C)+P(D)+P(E)\\
      &=P(A\cup B).
  \end{aligned}
$$
よって、${P(A\cup B)=P(A)+P(B)-P(A\cap B)}$ が成り立つ。
\clearpage
\endgroup

\handouttitle{第2回\quad 確率変数と確率分布}

\keybox{この回では、試行の結果 $\omega$ から観測値 $X(\omega)$ を取り出し、
その値がどの範囲に入るかを確率で表す。
確率変数 $X$、分布関数 $F_X$、確率分布 $\mu_X$ の役割を区別しよう。}

\topic{確率変数：試行の結果から観測値を取り出す}
以下、$(\Omega,\F,P)$ を確率空間とする。
$\Omega$ は試行の結果全体、$\F$ は確率を測ることのできる事象の集合族、
$P$ は各事象に確率を与える測度である。
写像 $X:\Omega\to\R$ は、一つの結果 $\omega$ に対して一つの実数 $X(\omega)$ を対応させる。
たとえば、コインを2回投げた結果 $(H,T)$ に、表の個数 $1$ を対応させる。

観測値が集合 $B\subseteq\R$ に入るという条件は、標本空間上の集合
$$
  \{X\in B\}:=X^{-1}(B)
  =\{\omega\in\Omega:X(\omega)\in B\}
$$
を表す。$X^{-1}(B)$ は $B$ の逆像である。この集合が $\F$ に属するとき、初めて $P(X\in B)$ を定義できる。

\topic{可測性：観測値に関する条件を事象として扱う}
$\mathcal B(\R)$ を、実数の開区間から生成される最小の $\sigma$-加法族
（ボレル集合族）とする。開区間だけでなく、閉区間、半開区間、一点集合なども含まれる。
写像 $X:\Omega\to\R$ が実数値確率変数であるとは、
$$
  X^{-1}(B)\in\F\qquad\bigl(\text{すべての }B\in\mathcal B(\R)\bigr)
$$
を満たすことである。この条件を $X$ の可測性という。
可測性があれば、$\{X\le x\}$、$\{a<X<b\}$、$\{X=x\}$ などはすべて事象になる。

\examplelabel{可測性を確認するための同値な条件：}
次のいずれかを確認すれば、$X$ は確率変数である。
$$
  \begin{aligned}
    &\{a<X<b\}\in\F\qquad(\text{すべての }a,b\in\R,\ a<b),\\
    &\{X\le x\}\in\F\qquad(\text{すべての }x\in\R).
  \end{aligned}
$$
逆像は補集合と可算和を保ち、開区間も半直線 $(-\infty,x]$ も
$\mathcal B(\R)$ を生成するためである。
特に、$\Omega$ が有限または可算で $\F=2^\Omega$ のときは、
どの写像 $X:\Omega\to\R$ も確率変数になる。
また、$X$ が確率変数で $g:\R\to\R$ が連続なら、
$$
  (g\circ X)^{-1}(B)=X^{-1}\bigl(g^{-1}(B)\bigr)\in\F
  \qquad\bigl(B\in\mathcal B(\R)\bigr)
$$
なので、$g(X)$ も確率変数である。

\topic{例：公平なコインを独立に2回投げる}
$H$ を表、$T$ を裏とし、
$$
  \Omega=\{H,T\}^2,\qquad \F=2^\Omega,\qquad
  P(\{\omega\})=\frac14\quad(\omega\in\Omega)
$$
とする。表の個数を表す確率変数 $X$ を次のように定義する。
\begin{center}
\begin{tabular}{@{}c c c c c@{}}
\toprule
試行の結果 $\omega$ & $(H,H)$ & $(H,T)$ & $(T,H)$ & $(T,T)$ \\
\midrule
観測値 $X(\omega)$ & $2$ & $1$ & $1$ & $0$ \\
\bottomrule
\end{tabular}
\end{center}
たとえば $\{X=1\}=\{(H,T),(T,H)\}$ なので、
$$
  P(X=1)=P\bigl(\{(H,T),(T,H)\}\bigr)=\frac14+\frac14=\frac12.
$$
同じ観測値を与える結果を集めて、その確率を足し合わせることが計算の基本である。

\newpage
\topic{分布関数}
確率変数 $X$ の分布関数を
$$
  F_X(x):=P(X\le x)
  =P\bigl(\{\omega\in\Omega:X(\omega)\le x\}\bigr)
  \qquad(x\in\R)
$$
と定める。$x$ は確定したしきい値であり、$F_X(x)$ は $0$ から $1$ の間の数である。
$F_X(x)$ は、$X$ の値が $x$ 以下となる確率を累積したものである。

\examplelabel{非減少性}
$$
  x\le y\quad\Longrightarrow\quad
  \{X\le x\}\subseteq\{X\le y\}
  \quad\Longrightarrow\quad F_X(x)\le F_X(y).
$$

\examplelabel{右連続性と無限遠での極限：}
証明には、確率測度の次の連続性を使う。
事象列 $(A_n)_{n\ge1}$ に対し、
$$
  \begin{aligned}
    A_n\uparrow A&\quad\Longrightarrow\quad P(A_n)\uparrow P(A),\\
    A_n\downarrow A&\quad\Longrightarrow\quad P(A_n)\downarrow P(A).
  \end{aligned}
$$
$A_n\uparrow A$ は $A_n\subseteq A_{n+1}$ かつ $\bigcup_n A_n=A$、
$A_n\downarrow A$ は $A_{n+1}\subseteq A_n$ かつ $\bigcap_n A_n=A$ を意味する。
確率測度では $P(\Omega)=1<\infty$ なので、減少列にも適用できる。
これらは可算加法性から従う。

任意の $x\in\R$ に対して
$$
  \{X\le x+1/n\}\downarrow\{X\le x\}
  \quad\Longrightarrow\quad
  F_X(x+1/n)\longrightarrow F_X(x).
$$
さらに、$0<h\le1/n$ なら $F_X(x)\le F_X(x+h)\le F_X(x+1/n)$ なので、
$$
  \lim_{h\downarrow0}F_X(x+h)=F_X(x).
$$
また、$X$ が実数値であることから
$$
  \begin{aligned}
    \{X\le n\}\uparrow\Omega
      &\quad\Longrightarrow\quad F_X(n)\longrightarrow1,\\
    \{X\le-n\}\downarrow\varnothing
      &\quad\Longrightarrow\quad F_X(-n)\longrightarrow0.
  \end{aligned}
$$
非減少性と合わせると、
$\lim_{x\to-\infty}F_X(x)=0$、$\lim_{x\to+\infty}F_X(x)=1$ を得る。

\topic{区間の確率と一点の確率}
$a<b$ のとき、$\{X\le b\}=\{X\le a\}\cup\{a<X\le b\}$ は互いに素な和なので、
$$
  P(a<X\le b)=F_X(b)-F_X(a).
$$
左極限 $F_X(x-):=\lim_{t\uparrow x}F_X(t)$ は非減少性により存在する。
$\{X\le x-1/n\}\uparrow\{X<x\}$ から
$$
  F_X(x-)=P(X<x),\qquad P(X=x)=F_X(x)-F_X(x-).
$$
したがって、分布関数の点 $x$ での跳びの大きさは $P(X=x)$ に等しい。
端点を含むかどうかに注意すると、
$$
  P(a\le X\le b)=F_X(b)-F_X(a-),\qquad
  P(a<X<b)=F_X(b-)-F_X(a).
$$

\topic{コインの例の分布関数}
前のページの $X$ に対し、$P(X=0)=1/4$、$P(X=1)=1/2$、$P(X=2)=1/4$ である。
\par\smallskip
\noindent\begin{minipage}[c]{0.40\linewidth}
$$
  F_X(x)=
  \begin{cases}
    0, & x<0,\\
    \dfrac14, & 0\le x<1,\\
    \dfrac34, & 1\le x<2,\\
    1, & 2\le x.
  \end{cases}
$$
\end{minipage}\hfill
\begin{minipage}[c]{0.57\linewidth}
\centering
\begin{tikzpicture}[x=2cm,y=2.8cm,font=\small]
  % Axes and guides; jump segments are not part of the graph.
  \draw[muted!35,densely dotted] (0,0.25)--(2.8,0.25);
  \draw[muted!35,densely dotted] (0,0.75)--(2.8,0.75);
  \draw[muted!35,densely dotted] (0,1)--(2.8,1);
  \draw[ink,->,line width=0.4pt] (-0.75,0)--(3.05,0)
    node[right] {$x$};
  \draw[ink,->,line width=0.4pt] (0,-0.05)--(0,1.2)
    node[above] {$F_X(x)$};
  \node[below left,ink] at (0,-0.03) {$0$};
  \foreach \x in {1,2}{
    \draw[ink] (\x,0)--(\x,-0.035);
    \node[below,ink] at (\x,-0.035) {$\x$};
  }
  \node[left,ink] at (-0.06,0.25) {$\frac14$};
  \node[left,ink] at (-0.06,0.75) {$\frac34$};
  \node[left,ink] at (-0.06,1) {$1$};
  % Four horizontal pieces, with arrows for the two unbounded tails.
  \draw[accent,<-,line width=1pt] (-0.7,0)--(0,0);
  \draw[accent,line width=1pt] (0,0.25)--(1,0.25);
  \draw[accent,line width=1pt] (1,0.75)--(2,0.75);
  \draw[accent,->,line width=1pt] (2,1)--(2.85,1);
  \foreach \x/\y in {0/0,1/0.25,2/0.75}{
    \draw[ink,fill=white,line width=0.7pt] (\x,\y) circle[radius=2pt];
  }
  \foreach \x/\y in {0/0.25,1/0.75,2/1}{
    \fill[ink] (\x,\y) circle[radius=2pt];
  }
\end{tikzpicture}
\end{minipage}
\par\smallskip
{\small\color{muted}黒丸はその点での値、白丸は含まない端点を表す。各跳びの大きさは $P(X=k)$（$k=0,1,2$）に等しい。}\par
たとえば $F_X(1)=P(X=0)+P(X=1)=3/4$ であり、$P(X=1)=1/2$ とは異なる。
また、$F_X(1)-F_X(1-)=3/4-1/4=1/2$ が点 $1$ での跳びである。

\newpage
\topic{確率分布：観測値の空間上の確率測度}
確率変数 $X$ の確率分布 $\mu_X$ を、$\mathcal B(\R)$ 上で
$$
  \mu_X(B):=P(X\in B)=P\bigl(X^{-1}(B)\bigr)
  \qquad\bigl(B\in\mathcal B(\R)\bigr)
$$
と定める。$P$ は標本空間 $\Omega$ の事象に確率を与え、
$\mu_X$ は観測値の空間 $\R$ のボレル集合に確率を与える。
標本空間上の確率を写像 $X$ によって移したもので、押し出し測度ともいう。

\examplelabel{$\mu_X$ が確率測度となる理由：}
可測性により右辺は定義でき、$\mu_X(B)\ge0$ である。また、
$$
  \mu_X(\R)=P\bigl(X^{-1}(\R)\bigr)=P(\Omega)=1.
$$
互いに素なボレル集合 $B_1,B_2,\ldots$ の逆像も互いに素なので、
$$
  \begin{aligned}
    \mu_X\!\left(\bigcup_{n=1}^{\infty}B_n\right)
      &=P\!\left(\bigcup_{n=1}^{\infty}X^{-1}(B_n)\right)\\
      &=\sum_{n=1}^{\infty}P\bigl(X^{-1}(B_n)\bigr)
       =\sum_{n=1}^{\infty}\mu_X(B_n).
  \end{aligned}
$$
したがって $(\R,\mathcal B(\R),\mu_X)$ も確率空間になる。

\topic{確率分布と分布関数の関係}
定義から
$$
  F_X(x)=\mu_X\bigl((-\infty,x]\bigr),\qquad
  \mu_X\bigl((a,b]\bigr)=F_X(b)-F_X(a)\quad(a<b).
$$
さらに、ボレル集合上の確率測度はこれらの区間上の値で一意に決まる
（測度の一意性定理）。したがって、分布関数は確率分布を一意に定める。

離散確率空間 $\Omega$（有限または可算、$\F=2^\Omega$）で
$p_\omega:=P(\{\omega\})$ とすると、
$$
  P(X=x)=\sum_{\omega\in\Omega:\,X(\omega)=x}p_\omega,
  \qquad
  F_X(t)=\sum_{\omega\in\Omega:\,X(\omega)\le t}p_\omega.
$$
コインの例では、任意の $B\in\mathcal B(\R)$ に対して
$$
  \mu_X(B)=\frac14\,\mathrm{1}_B(0)
           +\frac12\,\mathrm{1}_B(1)
           +\frac14\,\mathrm{1}_B(2),
$$
ただし $\mathrm{1}_B(x)$ は $x\in B$ のとき $1$、それ以外では $0$ をとる。

\topic{連続的な例：区間の長さから分布関数を求める}
$\Omega=[0,1]$、$\F=\mathcal B([0,1])$ とし、$P$ を $[0,1]$ 上の
ルベーグ測度 $\lambda$ とする。ここで $P(A)=\lambda(A)$ であり、
区間の確率はその長さに等しい。確率変数 $Z(\omega)=2\omega$ に対し、
$0\le z\le2$ なら
$$
  \{Z\le z\}=\{\omega\in[0,1]:2\omega\le z\}=[0,z/2],
  \qquad F_Z(z)=\lambda([0,z/2])=\frac z2.
$$
範囲外の $z$ も含めると
$$
  F_Z(z)=
  \begin{cases}
    0, & z<0,\\
    z/2, & 0\le z\le2,\\
    1, & z>2.
  \end{cases}
$$
これは $[0,2]$ 上の一様分布である。一点集合の長さは $0$ なので、
すべての $z\in\R$ に対して $P(Z=z)=0$ となる。
一方、区間の確率は $P(a<Z\le b)=(b-a)/2$（$0\le a<b\le2$）となる。

\keybox{計算の基本は、観測値の条件を $\omega$ の条件に書き換えることである。
離散モデルでは該当する標本の確率を足し、区間上の一様モデルでは該当する集合の長さを測る。}

\newpage
\topic{演習問題}
\begin{enumerate}[label=\arabic*.]
  \item 公平なサイコロを独立に2回投げ、出た目を $U,V$ とする。
  $X=U+V$、$Y=|U-V|$ に対し、
  分布関数 $F_X(x)$、$F_Y(x)$ を求めよ。
  \item $\Omega=[0,1]$、$\F=\mathcal B([0,1])$、$P=\lambda$（ルベーグ測度）とする。
  $X(\omega)=3\omega+1$、$Y(\omega)=(X(\omega)-5/2)^2$ の分布関数を求めよ。
  $Y$ の計算では、事象 $\{Y\le y\}$ を $\omega$ の区間として表せ。
\end{enumerate}

\topic{レポート問題}
\begin{enumerate}[label=\arabic*.]
  \item 公平なコインを独立に３回投げる。裏の数を表す確率変数を定義したい。確率空間およびその確率変数の定義を述べよ.
  \item 確率空間$(\Omega,\mathcal F,\mathbb P)$と写像 $X:\Omega\to \mathbb R$を考える. 次の集合族が$\sigma$加法族であることを示せ:
$$\{B\subset \mathbb R:~X^{-1}(B) \in \mathcal F\}.$$
\end{enumerate}

\newpage
\begingroup
\setlength{\parskip}{2.5pt}
\newcommand{\columnheading}[1]{\par\vspace{5pt}
  {\fontsize{11pt}{15pt}\selectfont\sffamily\mdseries\color{accent}#1}\par\vspace{2pt}}
\fancyhead[L]{\small\color{ink}\coursename{}・発展的内容}
\fancyhead[R]{\small\color{ink}第2回}

{\small\sffamily\mdseries\color{accent}発展的内容}\par
\vspace{4pt}
{\Large\color{ink}なぜ確率変数を「写像」として定義するのか}\par
\vspace{2pt}
{\color{accent}\rule{\linewidth}{0.7pt}}\par
\vspace{4pt}

大学3年生で初めて公理的確率論を学び、確率変数を「写像で定義する」と聞いたとき、
正直に言えば少し奇妙に感じた。入門的な本では $P(X=a)$ を直接与えることが多く、
標本空間からの写像を意識しなくても計算できる。そもそも写像という概念自体は、
確率に限ったものではない。それでも、確率変数を写像として定義すると、
試行の結果と、そこから取り出す観測値との関係がはっきりする。

\columnheading{観測としての「写像」}
日常語では「ランダムに数をとること」が確率変数だと説明される。
サイコロを1回投げるなら、$\Omega=\{1,\ldots,6\}$ として、
結果そのものを出目と見なせばよい。写像を明示するなら $X(\omega)=\omega$ である。

しかし2回投げるなら、$\Omega=\{1,\ldots,6\}^2$ とするのが自然だ。
結果 $(a,b)$ から、1回目の目、2回目の目、合計、差の絶対値を取り出す操作は、
それぞれ別の写像になる。
$$
  X_1(a,b)=a,\qquad X_2(a,b)=b,\qquad
  S(a,b)=a+b,\qquad D(a,b)=|a-b|.
$$
確率測度 $P$ は試行の結果に確率を与え、写像はその結果から観測値を取り出す。
この二つを組み合わせて、たとえば
$$
  P(S=3)=P\bigl(\{(a,b)\in\Omega:a+b=3\}\bigr)
$$
と計算する。$S$ 自体は固定した写像であり、同じ結果にはいつも同じ値を返す。
ランダムさは、どの結果が実現するかを $P$ で記述するところにある。

\columnheading{写像なら、複数の観測値を同時に扱える}
同じ確率空間上の確率変数 $X,Y$ があれば、和 $X+Y$、積 $XY$、
変換 $g(X)$ を各結果ごとに定義できる。
$g:\R\to\R$ がボレル可測なら $g(X)$ も確率変数である。
商 $X/Y$ を考える場合には、$Y=0$ での値の定め方も必要になる。

また、個々の分布だけでは、二つの観測値がどのように組み合わさるかは分からない。
公平なコイン1回の結果を $\Omega=\{H,T\}$ とし、
$$
  X(H)=1,\quad X(T)=0,\qquad Y=X,\qquad Z=1-X
$$
とおく。$W=X,Y,Z$ のいずれでも $P(W=0)=P(W=1)=1/2$ だが、
$$
  P(X=Y)=1,\qquad P(X=Z)=0.
$$
同じ標本空間上の写像として表せば、こうした同時の振る舞いも記述できる。
各変数の分布は $\mu_X(B)=P(X^{-1}(B))$ として、そこから取り出される。

\columnheading{可測性は「見えている情報」と結びつく}
コインを2回投げる例で、1回目の結果だけを知っているとする。
$A_H=\{(H,H),(H,T)\}$、$A_T=\{(T,H),(T,T)\}$ とおけば、
その情報で判定できる事象は
$$
  \F_1=\{\varnothing,\Omega,A_H,A_T\}
$$
である。1回目が表かどうかを表す $U(a,b)=\mathrm{1}_{\{H\}}(a)$ は
$\F_1$-可測だが、2回目が表かどうかを表す
$V(a,b)=\mathrm{1}_{\{H\}}(b)$ は $\F_1$-可測ではない。実際、
$$
  V^{-1}(\{1\})=\{(H,H),(T,H)\}\notin\F_1.
$$
$U$ は1回目の情報から値を決められるが、$V$ の値を知るには2回目の結果が必要だ。
一方、二回分の情報 $\F_2=2^\Omega$ に関しては、どちらも可測になる。
このように、可測性は「どの情報から観測値を決められるか」を表している。

\columnheading{おわりに}
初めは遠回りに見えるかもしれないが、具体例で計算していくうちに、
写像の定義の便利さが少しずつ見えてくる。
私自身も、使い続けることでその意味を理解するようになった。
確率過程や条件付き期待値を学ぶときにも、
「試行の結果」「観測値」「利用できる情報」を区別するこの見方が役に立つ。
\label{lastpage}
\clearpage
\endgroup
\end{document}
